\section{Names used in this paper}\label{app:names}

\Cref{tab:names} lists the names of learners, members, references, sets and terms that the paper uses, and \cref{tab:sumdefs,tab:sumdefs2} define the targets of the sum-structured suites other than \code{dlheld7}, which \cref{tab:dlheld7} defines. Every target is defined in \repo{study/boolfuncs.py}.

\begin{table}[p]
\centering\scriptsize\renewcommand{\arraystretch}{1.05}
\begin{tabularx}{\linewidth}{>{\raggedright\arraybackslash}p{3.6cm}>{\raggedright\arraybackslash}X}
\toprule
name & meaning\\
\midrule
\multicolumn{2}{l}{\emph{Learners}}\\
v1, v2, v3 & the study's portfolios of Part 8, 9 and 10: \code{hyb_portfolio}, \code{hyb_portfolio2}, \code{hyb_portfolio3} in the code (\cref{tab:members})\\
DFG & discrete feature growth, the Boolean member (\cref{sec:dfg}); v1.1 in v1, v1.2 in v2, v1.3 in v3\\
\code{p10_sum} & the Part 10 toy builder's candidate (DFG v1.3); v3's code on toy inputs\\
\code{p10_codes} & the adopted Part 10 image candidate (adds \code{cgbt}); v3's code on images\\
\code{p10_hier} & the other Part 10 image candidate (a max-pooling level in \code{ggbt}'s bank); not adopted\\
\code{p9_sum}, \code{p9_gates} & v2's Part 9 components on toy inputs and on images (v2's code there)\\
\addlinespace
\multicolumn{2}{l}{\emph{Members of v3} (\cref{tab:members})}\\
\code{dfg}, \code{dfg_forced} & DFG and its small-$n$ variant that forces the best candidate\\
\code{ggbt} & the blob booster: LightGBM on pixels, pooled sums and prototype pairs\\
\code{cgbt} & the bar booster: LightGBM on pixels, pooled sums and bar sums (new in v3)\\
\code{gknn_k} & nearest-neighbour votes on the 9-pixel pooled sums\\
\code{gbt}, \code{et}, \code{knn_<metric>_k}, \code{lin_C} & Part 8's LightGBM, ExtraTrees, nearest-neighbour votes and logistic regression on the raw columns\\
\addlinespace
\multicolumn{2}{l}{\emph{References} (\cref{sec:refs})}\\
\code{nets}, \code{nets_wd}, \code{nets_cewd}, \code{nets_cewd0.3}, \code{nets_cewd1.0} & the toy networks: square loss; weight decay 0.1; logistic loss with weight decay 0.1, 0.3, 1.0\\
\code{nets_best} & per target, the network variant with the smallest $\nstar$ among those run (hindsight)\\
\code{lin}, \code{hgb} & calibration references: cross-validated logistic regression on the inputs (a separate run, not the member \code{lin_C}); default histogram gradient boosting\\
\code{mlp}, \code{mlp_reg} & Part 7's baseline MLP (256-256) and its L2-regularized version\\
\code{mlp_strong}, \code{mlp_wide}, \code{mlp_xl} & ensembles of three MLPs at width 256, 512 and 1024; \code{mlp_strong1} and \code{mlp_wide1} are their single networks; \code{mlp_xl} is CPU-matched\\
\code{mlp_best}, \code{mlp_best+} & per-cell minimum over the Part 7--8 MLP variants (hindsight); \code{mlp_best+} adds the wider ones\\
\code{mlp_pool} & control: \code{mlp_strong} given v2's pooled sums\\
\code{cnn} & context: a small CNN on the unpermuted pixels\\
\addlinespace
\multicolumn{2}{l}{\emph{Sets} (\cref{sec:tasks,tab:status})}\\
\code{designed}, \code{random}, \code{dlheld}, \code{dlheld2}--\code{dlheld7}, \code{dlsymdev}, \code{dlfamdev} & the 11 toy suites; ``the blind family'' is the generator of \code{dlheld6}, \code{dlfamdev} and \code{dlheld7}\\
HELDOUT2; the Part 7 set & the 10 fresh tabular datasets of Part 8; the 11 held out in Part 7\\
\code{qd_dev}, \code{qd_a}--\code{qd_d}; Fashion & the five QuickDraw sets; Fashion-MNIST\\
target prefixes & f7\_ \code{dlheld7}; h6\_ \code{dlheld6}; h5 \code{dlheld5}; h4 (controls) \code{dlheld4}; fd\_ \code{dlfamdev}; h2rf, h3rf, hrf random formulas of \code{dlheld2}, \code{dlheld3}, \code{dlheld}; rf \code{random}. \code{dlsymdev}, \code{dlheld4} and the earlier suites use descriptive names (sym7, band11, tail13, majsym, thr10, maj9, mux6, \dots)\\
\addlinespace
\multicolumn{2}{l}{\emph{Terms}}\\
$\nstar$ & training-set size at which the seed-mean test error reaches 10\% (\cref{sec:metrics})\\
cell; clearly better, tie & a (set, $n$) pair; the pre-registered verdicts (\cref{sec:metrics})\\
budget & one toy network run, or three runs of the baseline \code{mlp}, in FLOPs (\cref{sec:metrics})\\
character; block & $\chi_S(x)=\prod_{i\in S}x_i$; indicators of the groups of a signed sum $t=\sum_{i\in S}s_ix_i$, folded ($|t|$), modulated, product, split or weighted (\cref{sec:dfg})\\
sum-structured target & one of the 54 targets of \code{dlsymdev}, \code{dlheld4}--\code{dlheld7} and \code{dlfamdev} that is not a control\\
$S_i(k)$; bar & the $k$-pixel pooling set of pixel $i$; a sum along a recovered line, half-length $L$ (\cref{sec:image})\\
$n_{\rm struct}$, $n_{\rm detect}$, share, reach & discovery size, detection size, a term's share and the search's reach (\cref{sec:interp-toy})\\
Part $k$ & the $k$-th part of the study; Parts 6--9 are reported in \cite{discretelearners2026}, Parts 1--5 in \cite{gdcircuit2026}\\
\bottomrule
\end{tabularx}
\caption{Names used in this paper.}
\label{tab:names}
\end{table}

\begin{table}[p]
\centering{\scriptsize\renewcommand{\arraystretch}{1.1}
\begin{tabularx}{\linewidth}{ll>{\raggedright\arraybackslash}X}
\toprule
suite & target & definition\\
\midrule
\code{dlsymdev} & sym7 & |x1 + ... + x7| $\le$ 1\\
 & ex3of8 & exactly 3 of x1..x8 are +1\\
 & band11 & |x1 + ... + x11| in \{3,\allowbreak{} 5\}\\
 & far9 & |x1 + ... + x9| $\ge$ 5\\
 & ssym9 & |x1 - x2 + x3 - x4 + x5 + x6 - x7 + x8 + x9| $\le$ 1\\
 & symxpar & [|x1+...+\allowbreak{}x6| $\le$ 2] xor (x7 x8)\\
 & twosym & [|x1+...+\allowbreak{}x5| $\le$ 1] == [|x6+...+\allowbreak{}x10| $\ge$ 3]\\
 & thrsum & sign(x1+...+\allowbreak{}x6 + 3 x7 x8 + 1/2)\\
\addlinespace[3pt]
\code{dlheld4} & sym11 & |x1 + ... + x11| $\le$ 1\\
 & ex4of9 & exactly 4 of x1..x9 are +1\\
 & tail13 & |x1 + ... + x13| $\ge$ 5\\
 & band10 & |x1 + ... + x10| = 2\\
 & ssym8 & x1 + x2 - x3 + x4 - x5 - x6 + x7 - x8 = 0\\
 & symxor & [|x1+...+\allowbreak{}x7| $\le$ 1] xor x8\\
 & symand & [|x1+...+\allowbreak{}x6| $\le$ 2] \& (x7 | x8)\\
 & sumpar & sign(x1+...+\allowbreak{}x5 + 2 x6 x7 + 2 x8 x9 x10)\\
 & majsym & MAJ([|x1+...+\allowbreak{}x5| $\le$ 1],\allowbreak{} x6 x7,\allowbreak{} x8)\\
 & dnf3x3 & (x1 \& $\neg$x2 \& x3) | (x4 \& x5 \& $\neg$x6) | ($\neg$x7 \& x8 \& x9) \emph{(control)}\\
 & h4rf0 & random 9-gate formula on 7 inputs \emph{(control)}\\
 & h4rf1 & random 9-gate formula on 7 inputs \emph{(control)}\\
\addlinespace[3pt]
\code{dlheld5} & h5s0 & [$\sum_{i\le12}s_ix_i$ in [$-$10,\allowbreak{} $-$6,\allowbreak{} 0,\allowbreak{} 2,\allowbreak{} 4,\allowbreak{} 6,\allowbreak{} 8,\allowbreak{} 12]],\allowbreak{} s = [1,\allowbreak{} $-$1,\allowbreak{} 1,\allowbreak{} 1,\allowbreak{} $-$1,\allowbreak{} 1,\allowbreak{} $-$1,\allowbreak{} 1,\allowbreak{} 1,\allowbreak{} 1,\allowbreak{} 1,\allowbreak{} $-$1]\\
 & h5s1 & [$\sum_{i\le9}s_ix_i$ in [$-$9,\allowbreak{} $-$5,\allowbreak{} $-$1,\allowbreak{} 3,\allowbreak{} 9]],\allowbreak{} s = [$-$1,\allowbreak{} 1,\allowbreak{} 1,\allowbreak{} 1,\allowbreak{} $-$1,\allowbreak{} 1,\allowbreak{} $-$1,\allowbreak{} $-$1,\allowbreak{} $-$1] xor chi(10,\allowbreak{} 11)\\
 & h5s2 & [$\sum_{i\le7}s_ix_i$ in [$-$5,\allowbreak{} 1,\allowbreak{} 7]],\allowbreak{} s = [1,\allowbreak{} $-$1,\allowbreak{} $-$1,\allowbreak{} 1,\allowbreak{} 1,\allowbreak{} $-$1,\allowbreak{} $-$1]\\
 & h5s3 & [$\sum_{i\le13}s_ix_i$ in [$-$7,\allowbreak{} $-$5,\allowbreak{} 1,\allowbreak{} 5,\allowbreak{} 9,\allowbreak{} 13]],\allowbreak{} s = [1,\allowbreak{} 1,\allowbreak{} 1,\allowbreak{} $-$1,\allowbreak{} 1,\allowbreak{} 1,\allowbreak{} 1,\allowbreak{} $-$1,\allowbreak{} 1,\allowbreak{} $-$1,\allowbreak{} 1,\allowbreak{} $-$1,\allowbreak{} 1]\\
 & h5s4 & [$\sum_{i\le5}s_ix_i$ in [$-$5,\allowbreak{} $-$1,\allowbreak{} 5]],\allowbreak{} s = [$-$1,\allowbreak{} $-$1,\allowbreak{} $-$1,\allowbreak{} 1,\allowbreak{} 1] or $\neg$x6\\
 & h5s5 & [$\sum_{i\le13}s_ix_i$ in [$-$11,\allowbreak{} $-$7,\allowbreak{} $-$5,\allowbreak{} $-$3,\allowbreak{} 1,\allowbreak{} 9,\allowbreak{} 11]],\allowbreak{} s = [1,\allowbreak{} 1,\allowbreak{} 1,\allowbreak{} 1,\allowbreak{} $-$1,\allowbreak{} $-$1,\allowbreak{} $-$1,\allowbreak{} $-$1,\allowbreak{} $-$1,\allowbreak{} 1,\allowbreak{} $-$1,\allowbreak{} $-$1,\allowbreak{} $-$1]\\
 & h5s6 & [$\sum_{i\le9}s_ix_i$ in [$-$9,\allowbreak{} $-$5,\allowbreak{} $-$1,\allowbreak{} 1,\allowbreak{} 5]],\allowbreak{} s = [1,\allowbreak{} $-$1,\allowbreak{} 1,\allowbreak{} $-$1,\allowbreak{} $-$1,\allowbreak{} 1,\allowbreak{} 1,\allowbreak{} 1,\allowbreak{} 1]\\
 & h5s7 & [$\sum_{i\le13}s_ix_i$ in [$-$11,\allowbreak{} $-$9,\allowbreak{} $-$5,\allowbreak{} $-$1,\allowbreak{} 1,\allowbreak{} 5,\allowbreak{} 7,\allowbreak{} 13]],\allowbreak{} s = [$-$1,\allowbreak{} $-$1,\allowbreak{} 1,\allowbreak{} $-$1,\allowbreak{} 1,\allowbreak{} 1,\allowbreak{} $-$1,\allowbreak{} 1,\allowbreak{} 1,\allowbreak{} $-$1,\allowbreak{} $-$1,\allowbreak{} 1,\allowbreak{} $-$1]\\
 & h5t0 & sign($\sum_{i\le7}s_ix_i$ + 2 chi(8,\allowbreak{} 9) + 1/2),\allowbreak{} s = [$-$1,\allowbreak{} 1,\allowbreak{} 1,\allowbreak{} 1,\allowbreak{} 1,\allowbreak{} $-$1,\allowbreak{} 1]\\
 & h5t1 & sign($\sum_{i\le5}s_ix_i$ + 2 chi(6,\allowbreak{} 7) + 1/2),\allowbreak{} s = [1,\allowbreak{} $-$1,\allowbreak{} 1,\allowbreak{} $-$1,\allowbreak{} 1]\\
 & h5rf0 & random 9-gate formula on 7 inputs \emph{(control)}\\
 & h5rf1 & random 6-gate formula on 7 inputs \emph{(control)}\\
\bottomrule
\end{tabularx}
}
\caption{The targets of \code{dlsymdev}, \code{dlheld4} and \code{dlheld5}, as defined in \repo{study/boolfuncs.py}; \code{dlheld6} and \code{dlfamdev} are in \cref{tab:sumdefs2}, \code{dlheld7} in \cref{tab:dlheld7}. Conventions as there: $[\cdot]$ is $\pm1$, \texttt{*} multiplies $\pm1$ values, \texttt{c ? a : b} is a multiplexer, $\neg$ negates a literal, and in \code{dlheld5} $s$ lists the signs $s_i$ and $[t\in A]$ the accept set $A$. ``(control)'': one of the 14 non-symmetric controls; the other targets are the sum-structured ones, together with the 9 of \code{dlheld7}.}
\label{tab:sumdefs}
\end{table}

\begin{table}[p]
\centering{\scriptsize\renewcommand{\arraystretch}{1.1}
\begin{tabularx}{\linewidth}{ll>{\raggedright\arraybackslash}X}
\toprule
suite & target & definition\\
\midrule
\code{dlheld6} & h6\_00 & [((s+13)/2) mod 3 != 1],\allowbreak{} s = x1+\allowbreak{}x2+\allowbreak{}x3+\allowbreak{}x4+\allowbreak{}x5+\allowbreak{}x6+\allowbreak{}x7+\allowbreak{}x8\allowbreak{}$-$x9\allowbreak{}$-$x10\allowbreak{}$-$x11\allowbreak{}$-$x12+\allowbreak{}x13\\
 & h6\_01 & [s in \{$-$9,\allowbreak{}$-$7,\allowbreak{}$-$5,\allowbreak{}$-$1,\allowbreak{}1,\allowbreak{}7,\allowbreak{}9\}],\allowbreak{} s = $-$x1+\allowbreak{}x2\allowbreak{}$-$x3+\allowbreak{}x4+\allowbreak{}x5+\allowbreak{}x6\allowbreak{}$-$x7\allowbreak{}$-$x8\allowbreak{}$-$x9\\
 & h6\_02 & [1 $\le$ s $\le$ 5] * x10 * x11,\allowbreak{} s = x1+...+\allowbreak{}x9\\
 & h6\_03 & OR(AND($-$x4,\allowbreak{}x5),\allowbreak{} AND($-$x2,\allowbreak{}$-$x6),\allowbreak{} AND(x1,\allowbreak{}$-$x3,\allowbreak{}x4)) \emph{(control)}\\
 & h6\_04 & OR($-$x1,\allowbreak{}AND(x2,\allowbreak{}AND($-$x3,\allowbreak{}OR(x4,\allowbreak{}x5)))) \emph{(control)}\\
 & h6\_05 & if $-$x6 $\to$ +1; elif $-$x2 $\to$ $-$1; elif x8 $\to$ $-$1; elif $-$x4 $\to$ $-$1; elif $-$x3 $\to$ +1; elif x5 $\to$ $-$1; elif AND($-$x7,\allowbreak{}$-$x1) $\to$ $-$1; else +1 \emph{(control)}\\
 & h6\_06 & [s in \{$-$6,\allowbreak{}$-$2,\allowbreak{}0,\allowbreak{}6\}] * x7 * x8,\allowbreak{} s = x1+...+\allowbreak{}x6\\
 & h6\_07 & $-$x11 ? [s1 in \{$-$2,\allowbreak{}4\}] : [s2 $\le$ 0],\allowbreak{} s1 = x1\allowbreak{}$-$x2\allowbreak{}$-$x3\allowbreak{}$-$x4\allowbreak{}$-$x5\allowbreak{}$-$x6,\allowbreak{} s2 = x7+\allowbreak{}x8+\allowbreak{}x9\allowbreak{}$-$x10\\
 & h6\_08 & x11 ? [|s1| $\le$ 2] : [((s2+6)/2) mod 3 == 0],\allowbreak{} s1 = x1+...+\allowbreak{}x4,\allowbreak{} s2 = $-$x5\allowbreak{}$-$x6+\allowbreak{}x7\allowbreak{}$-$x8+\allowbreak{}x9\allowbreak{}$-$x10\\
 & h6\_09 & XOR([$-$4 $\le$ s1 $\le$ 0],\allowbreak{} [3 $\le$ s2 $\le$ 5]),\allowbreak{} s1 = x1+...+\allowbreak{}x6,\allowbreak{} s2 = $-$x7\allowbreak{}$-$x8+\allowbreak{}x9\allowbreak{}$-$x10\allowbreak{}$-$x11\\
 & h6\_10 & [|u| > 1],\allowbreak{} u = s + 3*($-$x13),\allowbreak{} s = x1+...+\allowbreak{}x12\\
 & h6\_11 & [s == 1],\allowbreak{} s = x1+...+\allowbreak{}x7\\
\addlinespace[3pt]
\code{dlfamdev} & fd\_00 & if AND($-$x2,\allowbreak{}x4) $\to$ $-$1; elif $-$x3 $\to$ +1; elif AND($-$x6,\allowbreak{}$-$x1) $\to$ $-$1; elif AND($-$x8,\allowbreak{}x7) $\to$ $-$1; elif $-$x5 $\to$ +1; else $-$1 \emph{(control)}\\
 & fd\_01 & [s $\ge$ 0],\allowbreak{} s = x1+...+\allowbreak{}x6\\
 & fd\_02 & [|s| > 4] * x9 * x10,\allowbreak{} s = $-$x1+\allowbreak{}x2+\allowbreak{}x3+\allowbreak{}x4+\allowbreak{}x5+\allowbreak{}x6+\allowbreak{}x7\allowbreak{}$-$x8\\
 & fd\_03 & OR(AND($-$x1,\allowbreak{}$-$x2),\allowbreak{}AND(OR(x3,\allowbreak{}AND($-$x4,\allowbreak{}$-$x5)),\allowbreak{}$-$x6)) \emph{(control)}\\
 & fd\_04 & [|s| $\le$ 0] * x7 * x8 * x9,\allowbreak{} s = x1+...+\allowbreak{}x6\\
 & fd\_05 & AND([s1 $\le$ 0],\allowbreak{} [|s2| $\le$ 3]),\allowbreak{} s1 = x1+...+\allowbreak{}x6,\allowbreak{} s2 = x7+...+\allowbreak{}x13\\
 & fd\_06 & XOR([s1 $\ge$ 2],\allowbreak{} [$-$5 $\le$ s2 $\le$ 1]),\allowbreak{} s1 = $-$x1\allowbreak{}$-$x2+\allowbreak{}x3\allowbreak{}$-$x4,\allowbreak{} s2 = $-$x5+\allowbreak{}x6+\allowbreak{}x7+\allowbreak{}x8+\allowbreak{}x9\allowbreak{}$-$x10\allowbreak{}$-$x11\\
 & fd\_07 & OR(AND(x1,\allowbreak{}$-$x6,\allowbreak{}x7,\allowbreak{}$-$x8),\allowbreak{} AND(x2,\allowbreak{}x3,\allowbreak{}$-$x5),\allowbreak{} AND(x1,\allowbreak{}$-$x4,\allowbreak{}$-$x8)) \emph{(control)}\\
 & fd\_08 & MAJ([$-$5 $\le$ s $\le$ 9],\allowbreak{} $-$x12,\allowbreak{} $-$x13),\allowbreak{} s = x1+\allowbreak{}x2\allowbreak{}$-$x3+\allowbreak{}x4+\allowbreak{}x5+\allowbreak{}x6\allowbreak{}$-$x7+\allowbreak{}x8+\allowbreak{}x9+\allowbreak{}x10\allowbreak{}$-$x11\\
 & fd\_09 & [s in \{2,\allowbreak{}6\}],\allowbreak{} s = x1+...+\allowbreak{}x6\\
 & fd\_10 & [((s+10)/2) mod 3 == 0],\allowbreak{} s = x1+\allowbreak{}x2+\allowbreak{}x3+\allowbreak{}x4\allowbreak{}$-$x5\allowbreak{}$-$x6+\allowbreak{}x7+\allowbreak{}x8+\allowbreak{}x9+\allowbreak{}x10\\
 & fd\_11 & OR([|s| > 3],\allowbreak{} x10),\allowbreak{} s = x1+...+\allowbreak{}x9\\
\bottomrule
\end{tabularx}
}
\caption{The targets of \code{dlheld6} and \code{dlfamdev}, as defined in \repo{study/boolfuncs.py}. Conventions as in \cref{tab:sumdefs}.}
\label{tab:sumdefs2}
\end{table}

\clearpage
